% =====================================================================
%  Abbildung 5.x  --  Biquad (fi.tf2)
% =====================================================================
\begin{figure}
\centering
\begin{tikzpicture}
  \node[fsig](in)   at (0.0,-0.25) {$x[n]$};
  \node[fb2] (sub)  at (2.8, 0.0)  {$-$};
  \coordinate       (sp)  at (5.6, 0.0);
  \node[fb]  (firb) at (7.8, 0.0)  {\texttt{fir}$(b_0,b_1,b_2)$};
  \node[fsig](out)  at (10.2,0.0)  {$y[n]$};
  \node[fb]  (fira) at (2.8,-1.9)  {\texttt{fir}$(a_1,a_2)$};
  \node[fb]  (mem)  at (4.9,-1.9)  {\texttt{mem}};

  \draw[fw] (in.east) -- \pinB{sub};
  \draw[fplain] (sub.east) -- (sp);
  \node[fdot] at (sp) {};
  \node[flab, anchor=south] at (5.6,0.12) {$w[n]$};
  \draw[fw] (sp) -- (firb.west);
  \draw[fw] (firb.east) -- (out);

  % Rueckkopplungszweig, gespiegelt gezeichnet (Eingang rechts, Ausgang links)
  \draw[fw] (sp) -- (5.6,-1.9) -- (mem.east);
  \draw[fw] (mem.west) -- (fira.east);
  \coordinate (jog) at (1.95,-1.9);
  \draw[fw] (fira.west) -- (jog) |- \pinA{sub};

  \begin{scope}[on background layer]
    \node[fgrpin, draw=none, fill=none,
          fit=(sub)(fira)(mem)(sp)(jog)]         (recf) {};
    \node[fgrp,   fit=(in)(recf)(firb)]          (proc) {};
    \node[fgrpin, inner sep=0pt, fit=(recf)]     (rec)  {};
  \end{scope}
  \fgname{proc}{process = fi.tf2(\ldots)}
  \fgname{rec}{sub \textasciitilde{} fir($a_1,a_2$)}
\end{tikzpicture}
\caption[Faust-Blockdiagramm des Biquad-Filters]%
{Biquad-Filter (\texttt{biquad.dsp}, \texttt{fi.tf2}). Die
Bibliotheksfunktion ist eine sequenzielle Komposition aus einem
Rekursionsblock und einem FIR-Block:
\texttt{sub \textasciitilde{} fir(a1,a2) : fir(b0,b1,b2)}. Beide greifen auf
dasselbe Zwischensignal~$w$ zu, weshalb nur eine Verz\"ogerungskette
entsteht -- das ist die Direktform~II. Die Direktform~I der manuellen
Implementierung (Abb.~\ref{fig:biquad-formen}) lie\ss{}e sich so nicht
ausdr\"ucken, weil ihr Vorw\"artszweig \"uber das Eingangs- und ihr
R\"uckkopplungszweig \"uber das Ausgangssignal l\"auft und sie deshalb zwei
getrennte Ketten ben\"otigt.}
\label{fig:faust-biquad}
\end{figure}
